\documentclass[12pt]{article}

\usepackage{sbc-template}
\usepackage{graphicx,url}
\usepackage[utf8]{inputenc}
\usepackage[T1]{fontenc}
\usepackage[brazil]{babel}
\usepackage{amsmath}

\sloppy

\title{Seleção de Características com Algoritmo Genético para Previsão de Cancelamento de Clientes}

\author{Bredley Bauer\inst{1}, Gabriel Coelho\inst{1}, Gustavo Haiduk Pereira\inst{1}}

\address{Centro Universitário Católica de Santa Catarina\\
  Jaraguá do Sul -- SC -- Brasil
}

\begin{document}

\maketitle

\begin{abstract}
  This paper proposes genetic-algorithm-based feature selection for customer
  churn prediction using the Telco Customer Churn dataset. Logistic regression
  trained with all eligible attributes will be compared with the same
  classifier trained on a subset selected by evolutionary search. The fitness
  function combines cross-validated F1-score with a penalty for the number of
  selected attributes. The evaluation will consider predictive performance,
  classification errors, feature reduction, and computational cost.
\end{abstract}

\begin{resumo}
  Este trabalho propõe o uso de um algoritmo genético para selecionar atributos
  de um modelo de previsão de cancelamento de clientes. Com o conjunto de dados
  Telco Customer Churn, será comparada uma regressão logística treinada com
  todas as variáveis elegíveis à mesma técnica aplicada a um subconjunto
  selecionado por busca evolutiva. A função de aptidão combina o F1-score em
  validação cruzada com uma penalização pelo número de atributos. A avaliação
  considerará desempenho preditivo, erros de classificação, redução de
  variáveis e custo computacional.
\end{resumo}

\section{Introdução}

Em serviços de telecomunicações, a perda de clientes, conhecida como
\textit{customer churn}, pode afetar a receita e orientar decisões de retenção.
A partir de dados históricos sobre contratos, serviços e cobranças, modelos
de classificação supervisionada podem estimar a propensão ao cancelamento.

Características pouco informativas ou redundantes podem aumentar a
complexidade do classificador. A seleção de atributos busca subconjuntos
relevantes, cuja utilidade deve ser verificada experimentalmente
\cite{guyon2003feature}.

Este trabalho investiga o uso de um algoritmo genético para selecionar
características na previsão de cancelamento, comparando uma regressão logística
com todos os atributos elegíveis e outra treinada com os atributos selecionados.
O estudo utiliza o conjunto \textit{Telco Customer Churn}, disponibilizado no
Kaggle \cite{blastchartelco}. A aplicação é voltada a analistas e equipes
de retenção de empresas com serviços recorrentes, especialmente do setor de
telecomunicações.

\section{Fundamentação Teórica}

\subsection{Classificação Supervisionada}

Na aprendizagem supervisionada, o modelo é treinado com exemplos que possuem
características de entrada e uma resposta conhecida. Na previsão de
cancelamento, o problema possui duas classes: clientes que cancelaram o
serviço, definidos como classe positiva, e clientes que permaneceram, definidos
como classe negativa.

A regressão logística foi escolhida como classificador por fornecer previsões
para duas classes, permitir a interpretação dos coeficientes e apresentar custo
de treinamento relativamente baixo. Essa última característica é útil porque
a seleção genética exige treinar o classificador repetidas vezes com diferentes
subconjuntos. O treinamento utilizará a biblioteca \textit{scikit-learn}
\cite{pedregosa2011scikit}.

\subsection{Seleção de Características}

Na seleção de características do tipo \textit{wrapper}, cada subconjunto é
avaliado pelo desempenho do classificador treinado com ele
\cite{kohavi1997wrappers}. A remoção de atributos pode simplificar o modelo e
reduzir os dados exigidos para novas previsões, mas também pode prejudicar o
desempenho. A comparação com todas as variáveis elegíveis permite medir esse
efeito \cite{guyon2003feature}.

\subsection{Algoritmos Genéticos}

Algoritmos genéticos são métodos de busca que mantêm uma população de soluções
candidatas. Cada indivíduo recebe uma medida de aptidão e participa, conforme
seu desempenho, da formação das gerações seguintes. A seleção favorece
candidatos mais aptos; o cruzamento combina partes de soluções; e a mutação
introduz alterações aleatórias \cite{holland1975adaptation}.

O elitismo pode preservar uma boa solução entre gerações. Como a busca é
estocástica, o resultado pode variar entre execuções e não há garantia de
ótimo global.

\section{Metodologia Proposta}

\subsection{Base de Dados}

O estudo utilizará o \textit{Telco Customer Churn}, publicado no Kaggle por
\textit{blastchar} \cite{blastchartelco}. A base disponibilizada contém
7.043 registros e 21 colunas: um identificador de cliente
(\texttt{customerID}), a variável de resposta (\texttt{Churn}) e 19 atributos
sobre perfil, contrato, serviços e cobranças.

A variável \texttt{Churn} será codificada como 1 para \texttt{Yes} e 0 para
\texttt{No}; \texttt{customerID} ficará fora da modelagem. A exploração dos
dados abrangerá tipos das colunas, valores ausentes, duplicatas e distribuição
das classes, orientando os tratamentos necessários.

\subsection{Preparação dos Dados}

A coluna \texttt{TotalCharges} será convertida para formato numérico, pois
pode conter valores registrados como texto. Valores numéricos ausentes serão
tratados, quando presentes, por imputação da mediana; para atributos
categóricos, será utilizada a categoria mais frequente. As variáveis
categóricas passarão por codificação \textit{one-hot}, e as numéricas serão
padronizadas para a regressão logística.

Será utilizada uma divisão estratificada de 80\% dos registros para
desenvolvimento e 20\% para teste, com semente aleatória registrada. A imputação, a
codificação e a padronização serão ajustadas apenas nos dados de treinamento
de cada partição de validação. Isso evita que informações reservadas à
avaliação influenciem o modelo \cite{scikitlearnpitfalls}.

A seleção genética atuará sobre as 19 colunas originais, antes da codificação
\textit{one-hot}. Assim, a escolha de uma variável categórica incluirá todas
as colunas derivadas dessa variável. O mesmo procedimento de preparação será
aplicado aos dois modelos comparados.

\subsection{Seleção por Algoritmo Genético}

Cada cromossomo será um vetor binário com uma posição por atributo original. O
valor 1 indica inclusão da variável no treinamento, enquanto 0 indica sua
exclusão. Para $d$ atributos, existem $2^d-1$ subconjuntos não vazios. O
algoritmo impedirá soluções sem atributos selecionados.

A aptidão de cada cromossomo será calculada a partir do F1-score médio da
classe positiva em validação cruzada estratificada com três partições,
realizada somente no conjunto de desenvolvimento. Para favorecer soluções com
menos atributos, será aplicada uma penalização proporcional ao tamanho do
subconjunto:

\begin{equation}
  \operatorname{fitness}(S) =
  \overline{F_1}_{\mathrm{CV}}(S) - \lambda\frac{|S|}{d},
  \label{eq:fitness}
\end{equation}

\noindent em que $S$ é o subconjunto selecionado, $|S|$ é sua quantidade de
atributos, $d$ é o total de atributos elegíveis e $\lambda$ controla a
penalização. O valor inicial de $\lambda$ será $0{,}02$.

A configuração inicial prevê 24 indivíduos, 15 gerações, seleção por
torneio de três indivíduos, cruzamento com probabilidade de $0{,}8$ e elitismo
de um indivíduo. Na mutação, cada indivíduo terá probabilidade de $0{,}2$ de
ser modificado; nesse caso, cada bit será invertido com probabilidade de
$1/d$. Esses parâmetros limitam o custo da busca, pois cada avaliação exige
treinar o classificador em três partições de validação. Eventuais ajustes
utilizarão somente os dados de desenvolvimento. A execução registrará a
aptidão por geração, os atributos selecionados e o tempo de busca.

\subsection{Avaliação dos Modelos}

A comparação utilizará regressão logística nos dois cenários: com todas as
características elegíveis e com o subconjunto selecionado pelo algoritmo
genético. Ambos serão treinados com os dados de desenvolvimento e avaliados no
mesmo conjunto de teste, mantido fora da seleção de atributos.

O F1-score da classe de cancelamento será a métrica principal, pois combina
precisão e revocação. A medida é útil quando as classes têm frequências
distintas, situação em que a acurácia isolada pode ocultar erros na classe de
interesse. Também serão examinados precisão, revocação, acurácia, matriz de
confusão e, quando
aplicável, a área sob a curva ROC.

A análise distinguirá falsos negativos, clientes que cancelam sem serem
identificados, e falsos positivos, clientes classificados como propensos ao
cancelamento que permanecem. Além das métricas, serão comparados o número de
atributos e o tempo de execução da otimização. A separação do teste também limita o risco de estimativas excessivamente
otimistas após a busca genética.

\subsection{Demonstração de Uso}

A demonstração será feita por um script em Python que receberá os atributos
de um novo cliente e retornará a classificação produzida pelo modelo treinado.
O script utilizará o mesmo processamento aplicado no treinamento e, quando
disponível, apresentará também a probabilidade estimada de cancelamento.

\section{Considerações Finais}

A proposta permitirá examinar o efeito da seleção genética sobre a quantidade
de atributos e o desempenho de uma regressão logística na previsão de
cancelamentos. A comparação também levará em conta o
custo computacional das sucessivas avaliações feitas pelo algoritmo genético.

A análise experimental ainda será realizada. Entre as limitações previstas
estão a representatividade do conjunto de dados, a variabilidade da busca
genética e o possível sobreajuste à validação. As previsões também não
estabelecem relações de causa e efeito entre as características dos clientes
e o cancelamento.

\bibliographystyle{sbc}
\bibliography{referencias}

\end{document}
